% ===================================================================
%  नकसा नं. १३ — होटल। --- भोजनालय। कॉफीघर।
%
%  Traduction, faite en 2026, en bhojpouri d'aujourd'hui, sur l'ido
%  de texto/io. Regles de la colonne : voir 10-nakasa-01.tex. Une
%  seule numerotation. 22 blocs, 86 renvois, vingt-neuf paires a
%  retourner. DEUX notes d'auteur, la premiere en tete, la seconde
%  posee entre les blocs 8 et 9 ; l'appel du bloc 9 se compose en
%  parentheses simples, « (*) », et non en exposant.
%
%  LE SAINT-EMILION DE L'ALINEA 9 EST UN DELIT AU BIHAR DEPUIS 2016,
%  ET LA COLONNE L'ECRIT. La colonne gujaratie avait pose la loi du
%  releve a son tableau 15 : la viande de bœuf y est punie de la
%  reclusion a perpetuite depuis 2017, le mot existe et sert tous les
%  jours, la colonne l'ecrit sans note et sans detour — LA LANGUE A
%  LES MOTS ; C'EST LE FAIRE QUI EST INTERDIT, PAS LE NOMMER. Le
%  Bihar, cœur du pays bhojpouri, applique depuis avril 2016 une
%  prohibition totale de l'alcool : vendre ou boire y est une
%  infraction penale. Le Gujarat en a une depuis 1960. DEUX COLONNES
%  INDIENNES DU RELEVE, DEUX ETATS OU CE VERRE EST INTERDIT, ET LES
%  DEUX L'ECRIVENT — la seconde en le sachant, parce que la premiere
%  avait deja etabli la regle.
%
%  Le reste du menu ne demande rien : झींगा les crevettes, ओझरी les
%  tripes, मछरी le poisson — le Bihar est un pays de poisson de
%  fleuve —, मेमना के जाँघ le gigot, पालक les epinards, पनीर le
%  fromage. La colonne levantine avait fait la meme constatation au
%  meme alinea et en avait tire que la difficulte gujaratie n'etait
%  pas la francite du repas mais les interdits. Cette troisieme
%  lecture le confirme par l'autre bout : ici l'interdit n'est pas
%  dans la table, il est dans le verre.
%
%  ET « शतरंज » (78) EST LE MOT REVENU. Le releve tient maintenant les
%  trois etapes d'une seule route. La colonne gujaratie avait note que
%  « શતરંજ » etait parti d'Inde — chaturanga, les quatre corps d'une
%  armee. La colonne levantine avait montre l'etape du milieu : le
%  sanscrit devenu شطرنج en arabe par le persan, et de cette forme-la
%  viennent scacchi, echecs, chess. Le bhojpouri ecrit शतरंज — QUI EST
%  LA FORME ARABO-PERSANE, RENTREE EN INDE AVEC LE SULTANAT. Le mot
%  est donc revenu chez lui meconnaissable, et c'est sous sa forme
%  etrangere qu'on joue au jeu indien dans la plaine ou il est ne. Le
%  sanscrit चतुरङ्ग survit a cote, mais il ne nomme plus le jeu.
%
%  ENFIN LE TRICTRAC (81) A COUTE UN REFUS, ET LE MOT REFUSE AVAIT
%  VOYAGE LUI AUSSI. « पचीसी » etait la tentation : deux camps, des
%  des, on fait le tour du plateau. Mais le trictrac n'est pas le
%  pachisi, et le pachisi est justement l'autre jeu indien parti pour
%  l'Europe — il y est devenu le Ludo et le Parcheesi. On ecrit donc
%  « बैकगैमन ». On modernise la langue, jamais les choses, meme quand
%  la chose refusee a la meme histoire que celle qu'on ecrit.
% ===================================================================

%%K t13-noto-1 noto
\VUnotes{143.2mm}{%
\textit{\VUgras{(*) कमरा के बिस्तार खातिर नकसा नं. ६ देखीं।}}%
}

%%K t13-apar-1 apar
\VUsaut{3.0mm}
\VUtitre{15.0pt}{60}{नकसा नं. १३}
\VUsaut{1.6mm}
\VUpk{10.2pt}{40}{होटल। --- भोजनालय।}
\VUsaut{2.0mm}
\VUpk{10.2pt}{40}{कॉफीघर।}
\VUsaut{3.4mm}

%%K t13-01-1 p
१. --- कल्ह साँझ के रोयान पहुँच के हम « \textit{होटल दे ल'ओसेआँ} » में
ठहरनी, जेकर सलाह हमरा एगो दोस्त देले रहलें।

%%K t13-01-2 p
हमरा एगो सुंदर \VUgras{कमरा} \textsuperscript{(2)} मिलल, दुसरका
\VUgras{मंजिल} \textsuperscript{(1)} पर, जहाँ हम बहुते नीक से सुतनी (*)।

%%K t13-02-1 p
२. --- आज भोर, आपन कमरा से निकलत, जहाँ \VUgras{नौकरानी}
\textsuperscript{(3)} जलपान ले आइल रहली, हम \VUgras{सिगरेट पिए के कमरा}
\textsuperscript{(5)} में गइनी, जवना के \VUgras{पढ़े के कमरा}
\textsuperscript{(4)} के तौर पर भी इस्तेमाल कइल जाला, आ ओहिजा
\VUgras{अखबार} \textsuperscript{(6)} पढ़नी, \VUgras{सिगरेट}
\textsuperscript{(8)} पियत। पढ़ला के बाद हम \VUgras{लिफ्ट}
\textsuperscript{(9)} से नीचे उतरनी आ सहर घूमे निकल गइनी।

%%K t13-03-1 p
३. --- काहेकि हम होटल के \VUgras{मुनीम} \textsuperscript{(12)} से पूछल
भुला गइल रहनी कि \VUgras{साझा मेज} \textsuperscript{(11)} पर खाना कवना बेरा
मिलेला, हम जल्दी लवट अइनी आ एह मउका पर होटल के \VUgras{नहाए के कमरा}
\textsuperscript{(10)} में नहाए चल गइनी। ओकरा बाद हम फेरु \VUgras{गल्ला}
\textsuperscript{(15)} लगे गइनी ताकि आपन एगो दोस्त के फोन कर सकीं। बाकिर
\VUgras{टेलीफोन} \textsuperscript{(13)} बेस्त रहे ; हमार परेसानी देख के
\VUgras{गल्ला वाली} \textsuperscript{(14)}, जे एगो \VUgras{बिल}
\textsuperscript{(17)} \VUgras{जातरी के रजिस्टर} \textsuperscript{(16)}
पर लिखत रहली,
\VUgras{दरबान} \textsuperscript{(18)} से कहली कि \VUgras{छोटका नौकर}
\textsuperscript{(19)} के भेज दें जे हमार काम \VUgras{साइकिल से}
\textsuperscript{(20)} कर आवे।

%%K t13-04-1 p
४. --- हम ओकरा के धन्यवाद देनी आ बाहर निकलनी ताकि ओह \VUgras{कुली}
\textsuperscript{(24)} के बकसीस दे सकीं जे आपन \VUgras{हाथ के ठेला}
\textsuperscript{(25)} पर इस्टेसन से हमार \VUgras{टोपी के डिब्बा}
\textsuperscript{(26)}, \VUgras{हमार सूटकेस} \textsuperscript{(27)} आ
\VUgras{हमार जातरा के झोरा} \textsuperscript{(28)} ले आइल रहे।
\VUgras{डाकिया} \textsuperscript{(21)}, जे अभिए पहुँचल रहलें, हमरा एगो
\VUgras{चिट्ठी} \textsuperscript{(22)} देलें।

%%K t13-05-1 p
५. --- हम कुछ देर ले केवाड़ के देहरी पर, \VUgras{चिट्ठी के डिब्बा}
\textsuperscript{(23)} लगे, रुकल रहनी आ सड़क के आवाजाही देखत रहनी।
\VUgras{गाड़ीवान} \textsuperscript{(30)}, जे \VUgras{ओमनिबस}
\textsuperscript{(29)} के रहे, आपन गाड़ी पर एगो \VUgras{संदूक}
\textsuperscript{(31)} लादत रहे जवन एगो \VUgras{जातरी}
\textsuperscript{(32)} के रहे, जेकरा के \VUgras{होटल
के मालिक} \textsuperscript{(33)} विदा करे आइल रहलें। एगो जवान \VUgras{तार
वाला} \textsuperscript{(35)} आराम से होटल के एगो ग्राहक खातिर
\VUgras{तार} \textsuperscript{(34)} ले आइल ; आ एगो \VUgras{सैलानी}
\textsuperscript{(36)}, कान्ह पर आपन \VUgras{कैमरा} \textsuperscript{(38)}
लटकवले, आपन \VUgras{गाइड-पोथी} \textsuperscript{(37)} देखत रहे।

%%K t13-tit-1 sub
\VUsaut{4.0mm}
\VUpk{10.2pt}{40}{खाना के बेरा।}
\VUsaut{3.0mm}

%%K t13-06-1 p
६. --- जाली के सामने, एगो बेफिकिर \VUgras{जूता पालिस करे वाला}
\textsuperscript{(39)} आपन \VUgras{बोझ के हुक} \textsuperscript{(40)} आ आपन
\VUgras{पालिस के डिब्बा} \textsuperscript{(41)} के बीच झपकी लेत रहे, जबकि
एगो जवान \VUgras{धोबिन} \textsuperscript{(42)}, जे बियाढ़न के एगो
\VUgras{टोकरी} \textsuperscript{(43)} शायल रहली, ओह \VUgras{हाल के मुखिया}
\textsuperscript{(44)} के गंभीर सूरत पर हँसत रहली जे केवाड़ लगे ठाढ़
रहलें।

%%K t13-07-1 p
७. --- \VUgras{बावर्ची} \textsuperscript{(45)} के देख के, जे आपन
\VUgras{रसोई} \textsuperscript{(47)} से निकलत रहलें, जवन \VUgras{तहखाना}
\textsuperscript{(48)} में बा, ताकि \VUgras{रसोई के छोकड़ा}
\textsuperscript{(46)} के कवनो काम से भेज सकसु, हमरा इयाद आइल कि खाना के
बेरा नगीच बा, आ हम भोजनालय में घुसनी, ओह अँगना से होके जहाँ एगो
\VUgras{ड्राइवर} \textsuperscript{(50)} आपन \VUgras{मोटर} \textsuperscript{(49)}
लगववले रहे, जबकि एगो \VUgras{मिस्त्री} \textsuperscript{(52)},
\VUgras{साइकिल वाला} \textsuperscript{(53)} के बतावल के हिसाब से, ओह
\VUgras{तीन पहिया के मोटर} \textsuperscript{(51)} के बनावत रहे जेकर अभिए
दुर्घटना भइल रहे।

%%K t13-08-1 p
८. --- हम एगो जवान \VUgras{छोटका नौकर} \textsuperscript{(54)} से पूछनी,
जे \VUgras{बीयर के गिलास} \textsuperscript{(55)} लेके जात रहे, कि साझा मेज
बरामदा के नीचे बा कि ना, आ ओकर हामी सुन के हम एगो मेज पर जगहा लेनी आ
आपन खातिर बढ़िया खाना मँगवनी।

%%K t13-noto-2 noto
\VUnotes{143.2mm}{%
\textit{\VUgras{(*) सबदकोस में देहल पकवान के सूची के मदद से मास्टर
साहेब नीचे के खाना के सूची आसानी से बदल सकेलें।}}%
}

%%K t13-tit-2 sub
\VUsaut{4.0mm}
\VUpk{10.2pt}{40}{बढ़िया खाना।}
\VUsaut{3.0mm}

%%K t13-09-1 p
९. --- बैरा हमरा खाना के \VUgras{सूची} (*) आ \VUgras{सराब के सूची} ले
आइल। हम चुन के मँगवनी : सुरू में \VUgras{झींगा} आ \VUgras{मक्खन}, आ
पहिला पकवान में \VUgras{काँ के तरीका के ओझरी}। हम \VUgras{मछरी} ना
खइनी, बाकिर \VUgras{मेमना के जाँघ} के एगो हिस्सा माँगनी, आ
\VUgras{तरकारी} में \VUgras{मलाई वाला पालक}। ओकरा बाद, काहेकि
\VUgras{बीच के पकवान} में कवनो हमरा नीक ना लागल, हम तरह-तरह के मिठाई
मँगवनी — \VUgras{पनीर}, \VUgras{फर} आ \VUgras{बिस्कुट}। ऊपर से हम एगो
बढ़िया सैंत-एमिलियोँ \VUgras{सराब} भी लेनी, आ आपन खाना से बहुते खुस
होके हम मेज छोड़नी, \VUgras{बैरा} \textsuperscript{(57)} के नीक
\VUgras{बकसीस} \textsuperscript{(59)} देके।

%%K t13-10-1 p
१०. --- हमरा जाए खातिर तइयार देख के \VUgras{मैनेजर} \textsuperscript{(60)}
हमरा के सुझावलें कि बालकनी पर जाके \VUgras{महासागर} \textsuperscript{(62)}
के शानदार नजारा देख लीं। दूर मछरी मारे के \VUgras{नाव}
\textsuperscript{(63)}, \VUgras{पाल वाला जहाज} \textsuperscript{(64)} आ एगो
बड़हन \VUgras{भाप के जहाज} \textsuperscript{(65)} लउकत रहे, जेकर करिया
आकृति उजर \VUgras{बादर} \textsuperscript{(67)} पर उभरत रहे, जवन
\VUgras{छितिज} \textsuperscript{(66)} के रहे। हलुक हवा ओह \VUgras{झंडी}
\textsuperscript{(70)} के लहरावत रहे जवन \VUgras{तखती}
\textsuperscript{(69)} के ऊपर लागल रहे, आ ऊ तखती \VUgras{छाजन}
\textsuperscript{(68)} के रहे।

%%K t13-tit-3 sub
\VUsaut{3.0mm}
\VUpk{10.2pt}{40}{मन बहलाव।}
\VUsaut{3.0mm}

%%K t13-11-1 p
११. --- नजारा एतना मजेदार रहे कि हम तय कइनी कि काल्ह कॉफीघर के छत पर
\VUgras{दूरबीन} \textsuperscript{(71)} भा \VUgras{लमहर दूरबीन}
\textsuperscript{(72)} लेके चढ़ब।

%%K t13-12-1 p
१२. --- बालकनी छोड़ के हम होटल के कॉफीघर में गइनी ताकि एगो
\VUgras{पेय} \textsuperscript{(73)} पिये आ \VUgras{बिलियर्ड के खेल}
\textsuperscript{(75)} खेलीं। हम बिना रुकले ओह कॉफी-कमरा से गुजरनी जहाँ
बहुते ग्राहक \VUgras{शतरंज} \textsuperscript{(78)}, \VUgras{चेकर}
\textsuperscript{(79)}, \VUgras{डोमिनो} \textsuperscript{(80)},
\VUgras{बैकगैमन} \textsuperscript{(81)} भा \VUgras{ताश} \textsuperscript{(82)}
खेलत रहे, आ सोझे \VUgras{बिलियर्ड कमरा} \textsuperscript{(74)} में चढ़
गइनी।

%%K t13-13-1 p
१३. --- खेल खतम भइला पर हम भुइयाँ-तल पर उतरनी आ बैरा से चिट्ठी लिखे
खातिर \VUgras{लिफाफा} \textsuperscript{(84)}, \VUgras{कागज}
\textsuperscript{(83)}, \VUgras{टिकट} \textsuperscript{(85)}, \VUgras{कलम}
\textsuperscript{(87)} आ \VUgras{सियाही} \textsuperscript{(86)} माँगनी।

%%K t13-13-2 p
ओकरा बाद हम एगो \VUgras{कोचवान} \textsuperscript{(92)} के बोलवनी जेकर
\VUgras{गाड़ी} \textsuperscript{(88)} कॉफीघर के सामने खड़ा रहे। हम ओकरा से
घंटा के भा एक फेरा के भाड़ा पूछनी, आ जब भाड़ा तय हो गइल, ऊ हमरा खातिर
\VUgras{गाड़ी के केवाड़} \textsuperscript{(89)} खोललस आ हमरा के बइठवलस। ऊ
फेरु \VUgras{गद्दी} \textsuperscript{(91)} पर चढ़ल, \VUgras{चाबुक}
\textsuperscript{(93)} से घोड़न के तेज हँकलस आ हमनी के एगो मनमोहक घुमाई
पर निकल गइनी जा।
